%% kapa: exam=1 round=34 year=2023 period=2 form=A subject=appraisal_law subject_ko=감정평가관계법규 seq=2 number=2
%% source: Q-Net 감정평가사 기출문제 (cst00309) artlSeq=5211789 file=raw/1008/5211789/2023년도 제34회 감정평가사 1차시험 2교시.pdf
%% conversion: kordoc 4.24.51 -> tools/gen1.py (auto); review=proofread-sonnet; flags=none
\begin{question}{2}
국토의 계획 및 이용에 관한 법령상 광역도시계획에 관한 설명으로 옳은 것은?
\begin{choices}
  \item 광역계획권이 둘 이상의 시·도의 관할 구역에 걸쳐 있는 경우에는 관할 시·도지사가 공동으로 광역도시계획을 수립한다.
  \item 광역계획권을 지정한 날부터 2년이 지날 때까지 관할 시·도지사로부터 광역도시계획의 승인 신청이 없는 경우에는 국토교통부장관이 광역도시계획을 수립한다.
  \item 중앙행정기관의 장, 시·도지사, 시장 또는 군수는 국토교통부장관이나 도지사에게 광역계획권의 지정 또는 변경을 요청할 수 없다.
  \item 도지사가 시장 또는 군수의 요청에 의하여 관할 시장 또는 군수와 공동으로 광역도시계획을 수립하는 경우에는 국토교통부장관의 승인을 받아야 한다.
  \item 국토교통부장관, 시·도지사, 시장 또는 군수가 기초조사정보체계를 구축한 경우에는 등록된 정보의 현황을 3년마다 확인하고 변동사항을 반영하여야 한다.
\end{choices}
\end{question}
\answer{1}
